\documentclass{amsart}

% --- PACKAGES ---
\usepackage[margin=1in]{geometry}
\usepackage{amsmath, amsthm, amssymb}
\usepackage[colorlinks=true, allcolors=black]{hyperref}

% The following packages are included but not currently used.
% They can be enabled if figures are added in the future.
% \usepackage{graphicx}
% \usepackage{tikz}

% --- DOCUMENT METADATA ---
\hypersetup{
    pdftitle={Polylogarithm Values at a Golden Ratio-Based Argument},
    pdfauthor={Tristen Harr},
    pdfsubject={Complex Analysis, Special Functions, Number Theory},
    pdfkeywords={Infinite Series, Polylogarithm, Golden Ratio, Special Functions, Numerical Analysis, Transcendental Number Theory}
}

% --- THEOREM-LIKE ENVIRONMENTS ---
\theoremstyle{plain}
\newtheorem{theorem}{Theorem}[section]
\newtheorem{proposition}[theorem]{Proposition}
\newtheorem{lemma}[theorem]{Lemma}
\newtheorem{corollary}[theorem]{Corollary}
\newtheorem{conjecture}[theorem]{Conjecture}

\theoremstyle{definition}
\newtheorem{definition}[theorem]{Definition}

\theoremstyle{remark}
\newtheorem*{remark}{Remark}

% --- CUSTOM OPERATORS AND COMMANDS ---
\DeclareMathOperator{\Li}{Li} % Use DeclareMathOperator for standard function names
\newcommand{\LambdaG}{\Lambda_{G1}}
\newcommand{\phiG}{\phi}

% ======================================================================
% --- BEGIN DOCUMENT ---
% ======================================================================

\begin{document}

\title[Polylogarithms at a Golden Ratio Argument]{Polylogarithm Values at a Golden Ratio-Based Argument}
\author{Tristen Harr}
\date{Revised June 2025}

\begin{abstract}
This paper introduces and analyzes a specific complex constant, denoted $\LambdaG$, which is derived from fundamental properties of the golden ratio, $\phiG$. We define the constant's components based on inverse powers of $\phiG$, and we prove the key property that its magnitude is less than one. This validates the use of $\LambdaG$ as an argument in the Polylogarithm function, $\Li_s(z)$. We present high-precision numerical evaluations for the Dilogarithm ($s=2$) and Trilogarithm ($s=3$) of this constant. Based on this numerical evidence, we conjecture that the values $\Li_s(\LambdaG)$ are transcendental for all integers $s \ge 2$ and do not belong to the field extension $\mathbb{Q}(\pi, e, \ln(2), \phiG, \gamma)$. This line of inquiry is motivated in part by potential applications in the study of quasi-crystals.
\end{abstract}

\maketitle

% --- SECTION 1: INTRODUCTION ---
\section{Introduction}
The study of special functions and constants often reveals connections between different areas of mathematics. This paper focuses on the analysis of the Polylogarithm function at a novel complex argument, $\LambdaG$, derived from the golden ratio, $\phiG = (1+\sqrt{5})/2$. The motivation for this constant stems from the classical problem of dividing a line segment in extreme and mean ratio, as found in Book VI, Proposition 30 of Euclid's \textit{Elements}. We first rigorously define $\LambdaG$ from inverse powers of $\phiG$ and prove its validity as a polylogarithm argument. The analysis of the resulting values, $\Li_s(\LambdaG)$, forms the central topic of this work. Understanding the nature of these specific values may offer insights into mathematical domains where the golden ratio is foundational. For instance, given that the diffraction patterns of quasi-crystals are intrinsically linked to the golden ratio and that polylogarithms frequently arise in quantum field theory calculations (e.g., in the evaluation of Feynman diagrams), it remains an open and speculative question whether values like $\Li_s(\LambdaG)$ might emerge as fundamental components in the theoretical description of such aperiodic structures. The paper culminates in a formal conjecture about the transcendental nature of these values, based on high-precision numerical analysis.

% --- SECTION 2: THE CONSTANT AND ITS PROPERTIES ---
\section{The Constant \texorpdfstring{$\LambdaG$}{Lambda\_G1} and its Foundational Properties}

\subsection{Defining Components from the Golden Ratio}
We define two real constants, $T$ and $J$, directly from the algebraic properties of the golden ratio.
\begin{definition}[Primary Real Components]
The constants $T$ and $J$ are defined from the golden ratio, $\phiG$, as:
$$ T = \frac{1}{2\phiG} \quad \text{and} \quad J = \frac{1}{2\phiG^2} $$

\end{definition}

These definitions lead to a pair of linear relations that highlight their connection to classical geometry.
\begin{proposition}[System of Equations]
The constants $T$ and $J$ are the unique real numbers $(a,b)$ that satisfy the system:
\begin{align}
    a/b &= \phiG \\
    a+b &= 1/2
\end{align}

\end{proposition}
\begin{proof}
The ratio is immediate from the definitions: $T/J = (1/(2\phiG))/(1/(2\phiG^2)) = \phiG^2/\phiG = \phiG$. For the sum, we use the identity $\phiG+1 = \phiG^2$:
$$ T+J = \frac{1}{2\phiG} + \frac{1}{2\phiG^2} = \frac{\phiG+1}{2\phiG^2} = \frac{\phiG^2}{2\phiG^2} = \frac{1}{2} $$

The two equations $a/b = \phiG$ (or $a - \phiG b = 0$) and $a+b=1/2$ form a linear system in variables $(a,b)$. The determinant of the coefficient matrix is $\begin{vmatrix} 1 & -\phiG \\ 1 & 1 \end{vmatrix} = 1 - (-\phiG) = 1+\phiG$. Since $1+\phiG \neq 0$, the system has a unique solution.
\end{proof}

\begin{remark}
This formulation shows that the constraint $T+J=1/2$ is not an arbitrary choice, but a direct consequence of defining the components from successive inverse powers of $\phiG$. As discussed in Appendix A, this corresponds to the classical geometric problem of dividing a line segment of length $L=1/2$ in the golden ratio. The explicit values for $T$ and $J$ can be derived directly from their definitions:
$T = \frac{1}{2\phiG} = \frac{1}{1+\sqrt{5}} = \frac{\sqrt{5}-1}{(\sqrt{5}+1)(\sqrt{5}-1)} = \frac{\sqrt{5}-1}{4}$.
Using $\phiG^2 = \phiG+1 = \frac{3+\sqrt{5}}{2}$, we find $J = \frac{1}{2\phiG^2} = \frac{1}{3+\sqrt{5}} = \frac{3-\sqrt{5}}{(3+\sqrt{5})(3-\sqrt{5})} = \frac{3-\sqrt{5}}{4}$.
\end{remark}

From these constants, we define our primary object of study.
\begin{definition}[The Primary Constant]
The constant $\LambdaG$ is defined as the complex number $\LambdaG = T+iJ$.
\end{definition}

\subsection{Convergence Property}
A crucial property of $\LambdaG$ is that its magnitude is less than one, which is required for the convergence of the Polylogarithm series.
\begin{lemma}[Magnitude of $\LambdaG$]
\label{lem:convergence}
The constant $\LambdaG$ lies within the unit disk, i.e., $|\LambdaG| < 1$.
\end{lemma}
\begin{proof}
We require $|\LambdaG|^2 < 1$. We calculate the squared values of the components: $T^2 = \left(\frac{\sqrt{5}-1}{4}\right)^2 = \frac{6-2\sqrt{5}}{16}$ and $J^2 = \left(\frac{3-\sqrt{5}}{4}\right)^2 = \frac{14-6\sqrt{5}}{16}$. Summing these gives:
\begin{align*}
    |\LambdaG|^2 &= T^2+J^2 = \frac{6-2\sqrt{5} + 14-6\sqrt{5}}{16} \\
                 &= \frac{20-8\sqrt{5}}{16} = \frac{5-2\sqrt{5}}{4} \approx 0.131966
\end{align*}

Since $0.131966 < 1$, the condition is satisfied.
\qedhere
\end{proof}

% --- SECTION 3: DEFINING THE POLYLOGARITHM VALUES OF LAMBDA_G1 ---
\section{Defining the Polylogarithm Values of \texorpdfstring{$\LambdaG$}{Lambda\_G1}}
The property $|\LambdaG| < 1$ is central to this paper, as it allows us to define the Polylogarithm values at this constant. The Polylogarithm function is a key object in this study.
\begin{definition}[The Polylogarithm Function]
For any complex $s$, the Polylogarithm function of order $s$, $\Li_s(z)$, is defined by the power series:
$$ \Li_s(z) = \sum_{n=1}^{\infty} \frac{z^n}{n^s} $$

This series converges for all complex numbers $z$ such that $|z|<1$.
\end{definition}

Since $\LambdaG$ satisfies the convergence condition, we can define the specific numerical values that are the subject of our main conjecture:
$$ \Li_s(\LambdaG) = \sum_{n=1}^{\infty} \frac{(\LambdaG)^{n}}{n^s} $$
This identification holds for any complex s.

% --- SECTION 4: NUMERICAL ANALYSIS AND A CONJECTURE ---
\section{Numerical Analysis and a Conjecture on the Nature of \texorpdfstring{$\Li_s(\LambdaG)$}{Li\_s(Lambda\_G1)}}

\subsection{Numerical Values}
We present high-precision numerical evaluations for the first two non-trivial integer cases of the Polylogarithm of $\LambdaG$. The values were computed using the Python `mpmath` library with a working precision of 200 decimal digits to ensure robustness.
\begin{itemize}
    \item For $s=2$, the Dilogarithm of $\LambdaG$ is:
    $$ \Li_2(\LambdaG) \approx 0.322328253720993 + 0.226730769307749i $$
    
    \item For $s=3$, the Trilogarithm of $\LambdaG$ is:
    $$ \Li_3(\LambdaG) \approx 0.310345112613931 + 0.198884393433621i $$
    
\end{itemize}

\subsection{A Conjecture on the Nature of These Values}
The transcendental nature of special function values at algebraic arguments is a central topic in number theory. For context, values of the Dilogarithm are known to be transcendental at certain related real algebraic points, such as $\Li_2(1/2) = \frac{\pi^2}{12} - \frac{1}{2}\ln^2(2)$ and $\Li_2(\phiG^{-2}) = \frac{\pi^2}{15} - \ln^2(\phiG)$. The constant $\LambdaG$, being a non-real algebraic number that is not a root of unity, presents a distinct and challenging case. Transcendence proofs for polylogarithm values at real algebraic arguments often rely on techniques (such as specific functional equations or monodromy arguments) that do not readily extend to non-real complex arguments. This is largely because the rich structure of functional equations and the monodromy arguments they rely on become significantly more intricate in the complex plane away from the real line. This places the analysis of $\Li_s(\LambdaG)$ outside the scope of many established methods and helps motivate the present study.

To investigate the nature of these values, we employed the PSLQ integer relation algorithm, specifically using the implementation available in Python's `mpmath` library. A search for a null linear combination was performed on the high-precision real and imaginary parts of $\Li_s(\LambdaG)$ using a comprehensive basis of constants, including powers of $\pi, \ln(2),$ and $\phiG$, which are frequently found in known polylogarithm identities. No relation was found with a detection tolerance of $10^{-190}$. While not a formal proof, this failure provides strong numerical motivation for their apparent transcendental nature and leads us to propose the following general conjecture.
\begin{conjecture}
For any integer $s \ge 2$, the value $\Li_s(\LambdaG)$ is a transcendental number. Furthermore, it is conjectured that this number is not in the field $\mathbb{Q}(\pi, e, \ln(2), \phiG, \gamma)$.
\end{conjecture}
\begin{remark}
Proving this conjecture is beyond the scope of this paper, but proposing it frames a clear direction for future work.
\end{remark}

% --- SECTION 5: CONCLUSION ---
\section{Conclusion}
This paper has introduced a complex constant, $\LambdaG$, whose definition is derived non-arbitrarily from the golden ratio. We have proven that its magnitude is less than one, which validates its use as an argument in the Polylogarithm function. The primary contribution of this work is the introduction of this constant and the subsequent analysis of its polylogarithmic values, $\Li_s(\LambdaG)$. Based on high-precision numerical evidence from the $s=2$ and $s=3$ cases, we have proposed a new open problem: a conjecture on the transcendental nature of these values.

% --- SECTION 6: FUTURE WORK ---
\section{Future Work}
The conjecture presented in Section 4 opens a clear path for future investigation. A primary goal remains to prove or disprove the transcendental nature of $\Li_s(\LambdaG)$. More specific avenues for future work include:
\begin{itemize}
    \item Separately analyzing the real and imaginary parts of $\Li_s(\LambdaG)$. Key questions include whether these real values are themselves transcendental, or if they satisfy any currently unknown polynomial relationships with other well-known mathematical constants.
    \item Exploring other properties of the constant $\LambdaG$ itself.
    \item Investigating potential applications in physics, particularly in the study of quasi-crystals. Given that the golden ratio governs the diffraction patterns of these materials and that polylogarithms frequently arise in quantum field theory, exploring whether values like $\Li_s(\LambdaG)$ could emerge in theoretical models of such aperiodic structures may be a fruitful, albeit speculative, line of inquiry.
\end{itemize}

% --- APPENDICES AND REFERENCES ---
\appendix
\section{Geometric Motivation for the Defining Equations}
The definitions of $T$ and $J$ have a direct geometric motivation arising from the division of a line segment according to the golden ratio. If a segment of length $L$ is divided into two pieces, $a$ and $b$, such that $a/b = \phiG$, the pieces have lengths $a=L/\phiG$ and $b=L/\phiG^2$. Our constants $T$ and $J$ correspond to this construction for $L=1/2$. As shown in Proposition 2.2, the sum $T+J=1/2$ arises as a direct consequence of defining the components from the fundamental quantities $1/(2\phiG)$ and $1/(2\phiG^2)$. This provides a non-arbitrary basis for the constant's definition.

\section{Evaluation of a Related Polynomial-Weighted Series}
As a consequence of Lemma~\ref{lem:convergence}, $\LambdaG$ is a valid argument in the series whose sum is given by Eulerian polynomials, $A_k(x)$. This provides a direct evaluation for a family of series with $\LambdaG$ as the base. For any integer $k \ge 0$:
$$ \sum_{n=0}^{\infty} n^k \cdot (\LambdaG)^n = \frac{A_k(\LambdaG)}{(1-\LambdaG)^{k+1}} $$

While not central to the paper's main conjecture, this identity is an interesting property of the constant $\LambdaG$ and demonstrates its utility in evaluating other classes of well-known series.

\begin{thebibliography}{9}

\bibitem{Baker}
A. Baker. \textit{Transcendental Number Theory}. Cambridge University Press, 1975. 

\bibitem{Euclid}
Euclid, et al. \textit{The Thirteen Books of Euclid's Elements}. In \textit{Great Books of the Western World}, vol. 11, edited by R. M. Hutchins, Encyclopaedia Britannica, 1952. 

\bibitem{Lewin}
L. Lewin. \textit{Polylogarithms and Associated Functions}. North-Holland, 1981. 

\end{thebibliography}

\end{document}